%\documentclass[man, noextraspace]{apa6} 
\documentclass[man,floatsintext,noextraspace]{apa6}
\usepackage[natbibapa]{apacite} % Para manejar citas en estilo APA
\shorttitle{Thermo-Mechanical description of a Wankel Rotary Engine}
\title{Integrated Thermo-Chemical-Mechanical Modeling of a Four-Stroke Wankel Rotary Engine}
%\title{Thermo-Mechanical Description of a Wankel Rotary Engine: Computational simulation of the thermodynamic cycle and combustion}

\author{
Juan José Rojas Barón \\
Camilo Andres Bayona Roa \\
Sebastián Lemus Figueroa \\
Luis Alejandro Garzón Olaya 
}
\affiliation{Pontificia Universidad Javeriana}
\abstract{This study proposes an integrative methodology to simulate the thermo- chemical- mechanical behavior of a four-stroke Wankel rotary engine by combining epitrochoidal geometric modeling, finite-rate chemical kinetics, thermodynamic analysis, and a vectorial mechanics approach. This framework provides accurate predictions of instantaneous torque through vector cross-products between applied combustion forces and the effective rotor arm. Computational analysis of three distinct geometries demonstrates that eccentricity ($e$) is the primary driver of compression ratio and peak pressure, with the G3 configuration achieving a 52.8\% increase in $P_{max}$ compared to the baseline. Multi-fuel simulations indicate that methane offers superior thermodynamic efficiency and the lowest specific fuel consumption (0.031 g/kJ), while bioethanol shows a significant performance penalty due to its lower energy density. Validation against the AIE 650S commercial engine confirms that the model correctly captures power-RPM proportionality, though it highlights the necessity of incorporating high-speed limiting mechanisms, such as mechanical losses and reduced volumetric efficiency, for greater accuracy. This framework serves as a robust tool for optimizing design variables and evaluating performance under varying operational conditions.}



% Elimina el "Running head:" de la primera página
\makeatletter
\renewcommand{\ps@titlepage}{
  \renewcommand{\@oddhead}{} % Encabezado vacío
  \renewcommand{\@evenhead}{}
  \renewcommand{\@oddfoot}{\thepage} % Pie de página solo con número
  \renewcommand{\@evenfoot}{\thepage}
}
\makeatother
\usepackage{amsmath}
\usepackage{float}  
\usepackage{graphicx}
\graphicspath{{Images/}}


\begin{document}
\maketitle

\textbf{Keywords:} Wankel engine, thermodynamic cycle, vectorial mechanics, finite-rate chemical kinetics, two zone combustion, alternative fuels.

\section{1. Introduction}

The energy transformation sector faces significant challenges due to its reliance on conventional technologies, which often exhibit low thermal efficiencies, according to \cite{Wang2021}, and high levels of pollutant emissions. Traditional engines, such as Otto and Diesel cycle engines, frequently fail to optimize the thermal energy generated during the combustion process, leading to substantial energy losses and a considerable environmental footprint.

The use of unconventional technologies, such as Wankel rotary engines, presents an opportunity to rethink energy transformation processes. However, these engines remain underutilized, primarily due to historical concerns regarding their efficiency and durability. Furthermore, detailed analyses that consider both the thermodynamic and mechanical aspects of this technology under real-world conditions are scarce.

In the current context, where maximizing energy efficiency and minimizing environmental impact are priorities, it is essential to assess the potential of the Wankel engine as a viable alternative. Such an analysis would not only deepen the understanding of its behavior under various operational conditions but also provide a robust foundation for its optimization in practical applications, contributing to a more efficient and sustainable energy transformation paradigm.

This work proposes a comprehensive methodology that integrates geometric modeling, thermodynamic cycle simulation, finite-rate chemical kinetics, and vectorial mechanics. The thermo-chemical branch is initially developed through a homogeneous zero-dimensional One-Zone formulation coupled to the instantaneous chamber volume. This formulation constitutes the current thermo-chemical reference model and provides the baseline for the subsequent development of a Two-Zone combustion model and its final extension to a gasoline-lubricating-oil mixture. One of the main contributions lies in the incorporation of torque analysis based on vector cross-products between the applied force and the effective arm generated by rotor motion. By combining analytical functions—such as the epitrochoidal volume variation—with a dynamic mechanical model, this article offers a deeper and more accurate understanding of the power transmission mechanisms in Wankel engines. This integration enables the prediction of instantaneous torque, evaluation of performance under transient conditions, and ultimately supports the optimization of design variables with greater precision.


\subsection{1.1. State of the Art}

The Wankel engine is characterized by its innovative design, utilizing a rotor with curved sides that rotates eccentrically within an epitrochoidal-shaped housing. This motion generates three working chambers that alternately perform the four stages of the combustion cycle: intake, compression, combustion, and exhaust. One of its main advantages is its reduced number of moving parts compared to traditional piston engines, which significantly decreases friction and wear \citep{froede1961}. Moreover, it operates more smoothly with fewer vibrations due to the absence of reciprocating movements \citep{bensinger1973}. Additionally, its compact and lightweight design allows for greater power density \citep{froede1961}, and its ability to reach higher RPMs with ease results in a more consistent power output and an exhilarating driving experience \citep{yamamoto1981}. The simplicity of its mechanical design also contributes to lower maintenance requirements \citep{norbye1971}, making it particularly suitable for performance-oriented applications \citep{ludvigsen1973}. Finally, its continuous combustion cycle enables efficient power delivery, enhancing overall engine responsiveness and smooth acceleration \citep{kamo1986}.

\includegraphics{WE}

Academic investigations have focused on enhancing the performance and efficiency of Wankel rotary engines through various innovative approaches. Among these, analytical methods seek theoretical understanding, but the complex kinematics and thermodynamics of the engine have limited the development of general mathematical models. Experimental methods provide valuable operational insights but are costly and restricted to specific geometries. Computational models, meanwhile, enable detailed internal simulations and optimizations, though they are resource-intensive.

For instance, \cite{Shen2024} employed CFD to analyze combustion in Wankel engines and proposed turbulent jet ignition to mitigate incomplete combustion. \cite{Yang2022} explored hydrogen-enriched hybrid cycles, noting enhanced performance at high altitudes. \cite{Kim2024} applied 1D CFD in GT-POWER to optimize port design and fuel efficiency in a novel rotary engine derived from a gerotor pump.

Other studies include \cite{mukherjee2022}, who quantified intake losses in steam Wankel expanders and their impact on power output, and \cite{chen2020}, who built control-oriented models for the 225CS engine using hybrid analytical-neural approaches. \cite{bailey2018} analyzed direct injection limitations and fuel economy improvements in Wankel engines using one-dimensional modeling.

Technological advances continue with Mazda reintroducing rotary engines in 2023, LiquidPiston's XTS-210 model targeting military applications, and AIE and Freedom Motors refining air-cooled designs and low-emission engines. These efforts demonstrate renewed interest in rotary technology, emphasizing compactness, efficiency, and cleaner fuel options.

Despite these contributions, most research remains fragmented, focusing on isolated subsystems. There remains a need for integrative studies addressing geometry, combustion chemistry, and thermodynamic cycles within a unified modeling framework.The present work addresses this need progressively, beginning with a finite-rate One-Zone thermo-chemical formulation and establishing a common computational
framework that can subsequently be extended toward spatially differentiated Two-Zone combustion and gasoline--lubricating-oil mixtures.


\section{2. Methodology}

This section details the mathematical and physical methods used to describe, simulate, and analyze the behavior of a four-stroke Wankel rotary engine. The methodology is structured into four main components: (1) geometric formulation, (2) thermodynamic cycle modeling, (3) thermo-chemical combustion modeling using finite-rate chemical kinetics, progressing from a homogeneous One-Zone formulation to a Two-Zone burned/unburned formulation and, subsequently, to the incorporation of a gasoline-lubricating-oil mixture, and (4) torque and power analysis using vectorial mechanics. Together, these components provide a comprehensive framework to understand the performance, efficiency, and dynamic response of the engine under varying operating conditions.

\subsection{2.1 Geometric Formulation}

The geometric formulation of the Wankel engine is governed by a set of fundamental parameters that define the rotor–stator kinematic relationship and the instantaneous chamber volume. These parameters determine the architecture of the mechanism, the shape of the epitrochoidal housing, and the volumetric evolution throughout the thermodynamic cycle. In particular, they directly influence the maximum and minimum chamber volumes, and therefore the geometric compression ratio of the engine. A rigorous understanding of their geometric meaning and functional role is essential for the development of an accurate thermo–mechanical model and for properly characterizing the engine’s performance. 

The volumetric behavior and the resulting geometric compression ratio are primarily determined by four interconnected variables. First, the \textbf{eccentricity ($e$)} represents the distance between the rotor center and the eccentric shaft center, defining the orbital motion radius. An increase in $e$ amplifies the chamber volume variation and directly elevates the geometric compression ratio; therefore, its selection is critical to balance torque production and mechanical stresses. Complementing this, the \textbf{generating radius ($R$)} defines the rotor geometry and the formation of the stator epitrochoid. The $R/e$ ratio is a fundamental design constraint that ensures geometric compatibility while controlling the base chamber volume.


Furthermore, the \textbf{axial width ($W$)} of the rotor serves as a linear scaling factor that converts the instantaneous cross-sectional area into volume, directly influencing engine displacement without altering the compression ratio. Finally, the \textbf{parametric variable ($v$)} governs the mathematical definition of the housing's epitrochoidal profile, ensuring geometric consistency throughout the rotation. A rigorous integration of these variables allows for the precise characterization of the engine’s performance under varying configurations.

\subsubsection{2.1.1 Epitrochoidal Curve of the Stator}

The internal shape of the stator is defined by an epitrochoidal curve, governed by the following parametric equations \cite{yamamoto1981}:
\begin{align}
x &= e \cdot \cos(n\theta) + R \cdot \cos(\theta) \\
y &= e \cdot \sin(n\theta) + R \cdot \sin(\theta)
\end{align}

where \( e \) is the eccentricity, \( R \) is the radius of the base circle, and \( \theta \) is the rotation angle. This curve defines the trajectory of the rotor center and the envelope of the rotor's movement.

\subsubsection{2.1.2 Family of Curves Composing the Rotor}

The rotor’s profile ensures continuous contact with the epitrochoidal stator. Its geometry is described by the following equations \cite{yamamoto1981}:

\begin{align}
X &= R \cos(2v) + \frac{3}{2} \frac{e^2}{R} (\cos(8v) - \cos(4v)) 
+ e \sqrt{1 - \frac{9e^2}{R^2} \sin^2(3v)} \cdot (\cos(5v) + \cos(v)) \\
Y &= R \sin(2v) + \frac{3}{2} \frac{e^2}{R} (\cos(8v) + \sin(4v)) 
+ e \sqrt{1 - \frac{9e^2}{R^2} \sin^2(3v)} \cdot (\sin(5v) - \sin(v))
\end{align}

These equations are key to ensuring geometric compatibility and proper sealing during rotation.

\subsubsection{2.1.3 Volume Functions of the Chambers}

The instantaneous volume of each chamber varies with the rotor angle. The three chamber volumes are given by:

\begin{align}
V_1 &= \left[ \left(e^2 + \frac{R^2}{3}\right)\pi - \frac{\sqrt{3}}{4} R^2 \right] W 
      - \frac{3\sqrt{3}}{2} WeR \sin\left(2\left(\theta - \frac{2\pi}{3}\right) + \frac{\pi}{6}\right) \\
V_2 &= \left[ \left(e^2 + \frac{R^2}{3}\right)\pi - \frac{\sqrt{3}}{4} R^2 \right] W 
      - \frac{3\sqrt{3}}{2} WeR \sin\left(2\theta + \frac{\pi}{6}\right) \\
V_3 &= \left[ \left(e^2 + \frac{R^2}{3}\right)\pi - \frac{\sqrt{3}}{4} R^2 \right] W 
      - \frac{3\sqrt{3}}{2} WeR \sin\left(2\left(\theta + \frac{2\pi}{3}\right) + \frac{\pi}{6}\right)
\end{align}

These functions form the basis for the thermodynamic analysis of the Otto cycle.

\subsection{2.2 Baseline Ideal Thermodynamic Model}
The idealized Otto-type formulation is retained as a baseline model for comparison with the detailed thermo-chemical formulation introduced in the following section. In this baseline approach, compression and expansion are represented through polytropic relations and combustion is approximated as constant-volume heat addition. Therefore, this formulation should be interpreted as a simplified thermodynamic reference rather than as the final combustion model developed in this work.

\subsubsection{2.2.1 Polytropic Compression and Expansion}

Using a polytropic relation for adiabatic processes:

\begin{equation}
    P_2 = P_1 \left( \frac{V_1}{V_2} \right)^\gamma
\end{equation}

this equation is applied during the compression and expansion phases to estimate pressure evolution based on volume changes defined in the previous subsection.

\subsubsection{2.2.2 Heat Addition (Isochoric Process)}
Assuming constant volume during combustion, the final temperature is calculated using:
\begin{equation}
T_3 = \frac{Q}{m_{\text{total}} \cdot C_v} + T_2
\end{equation}
and the corresponding pressure is derived from:
\begin{equation}
P_3 = P_2 \cdot \frac{T_3}{T_2}
\end{equation}
The composition of the mixture, based on the air-fuel ratio (AFR), determines the molar mass and heat capacity of the reacting gases.


\subsection{2.3 Finite-Rate Thermo-Chemical Modeling Framework}

The thermo-chemical formulation extends the baseline ideal-cycle model by replacing prescribed instantaneous heat release with the explicit solution of finite-rate chemical kinetics. This approach allows the evolution of reactants, products, and intermediate species to be determined directly from the chemical mechanism instead of imposing a predefined burned-fraction law. Reduced and detailed kinetic mechanisms have been extensively employed in internal combustion engine simulations to describe fuel oxidation, ignition, and species formation under engine-relevant thermodynamic conditions
\cite{RaReitz2008,Mehl2009}.

The present stage of the model uses a homogeneous One-Zone formulation as the first level of thermo-chemical complexity. Pressure, temperature, and chemical composition are therefore assumed spatially uniform inside the modeled chamber at each eccentric-shaft angle. Single-Zone formulations provide a computationally efficient reduced-order description and have also been applied to rotary-engine
thermodynamic modeling \cite{Yang2022}.

This One-Zone formulation is retained as the thermo-chemical reference
configuration for the subsequent development of a Two-Zone model, in which burned and unburned regions will be resolved separately. The final planned extension will introduce a gasoline--lubricating-oil mixture within that Two-Zone framework.

\subsubsection{2.3.1 One-Zone Model}

\paragraph{Reactor formulation.}

The current thermo-chemical baseline represents the Wankel combustion chamber as a homogeneous zero-dimensional ideal-gas reactor. Under the One-Zone assumption, pressure, temperature, and chemical composition are spatially uniform within the complete chamber at each eccentric-shaft angle. The thermodynamic and chemical state is therefore described by

\begin{equation}
P=P(\theta_e), \qquad
T=T(\theta_e), \qquad
Y_k=Y_k(\theta_e),
\end{equation}

where $\theta_e$ is the eccentric-shaft angle and $Y_k$ is the mass fraction of chemical species $k$.

The reacting system is implemented using Cantera, an open-source software package for chemical kinetics, thermodynamics, and transport processes \cite{Cantera2025}. The reactor volume is imposed directly from the Wankel geometric model,

\begin{equation}
V=V(\theta_e),
\end{equation}

while time and eccentric-shaft angle are related through

\begin{equation}
\theta_e=\omega_e t,
\end{equation}

where $\omega_e$ is the eccentric-shaft angular velocity. This coupling ensures that chamber geometry, chemical kinetics, thermodynamic state, and indicated work are evaluated using the same angular reference.

\paragraph{Finite-rate chemical kinetics.}

Combustion is resolved through finite-rate chemical kinetics rather than through an imposed heat-release or burned-fraction function. For each chemical species $k$, the kinetic mechanism determines the net molar production rate $\dot{\omega}_k$. Expressed with respect to eccentric-shaft angle, the chemical source term is written as

\begin{equation}
\frac{dm_k}{d\theta_e}
=
\frac{\dot{\omega}_k W_k V}{\omega_e},
\end{equation}

where $m_k$ is the mass of species $k$, $W_k$ is its molecular weight,
$V$ is the instantaneous chamber volume, and $\omega_e$ is the
eccentric-shaft angular velocity.

Finite-rate formulations make it possible to follow the consumption and formation of individual chemical species throughout combustion, including intermediate species and radicals that cannot be represented by a prescribed global heat-release law \cite{McAllister2011,RaReitz2008}.

\paragraph{Reference fuel and kinetic mechanism.}

Commercial gasoline is a complex multicomponent fuel whose composition may include a large number of hydrocarbons with different chemical structures. Consequently, combustion simulations commonly employ surrogate fuels composed of a reduced number of well-characterized hydrocarbons in order to reproduce selected physical and chemical characteristics of the real fuel \cite{RaReitz2011,Daly2018}.

In the present One-Zone stage, iso-octane ($\mathrm{IC_8H_{18}}$) is adopted as a controlled reference hydrocarbon. Iso-octane is one of the primary reference fuel components commonly used in engine-combustion kinetic studies \cite{RaReitz2008}. It is not intended here
to reproduce the complete composition or combustion behavior of commercial gasoline. Instead, it provides a chemically defined baseline from which the thermo-chemical architecture can be developed and numerically verified before introducing a multicomponent gasoline surrogate and lubricating oil.

The working kinetic mechanism used in the present model contains 56 chemical species and 332 reactions. These values correspond specifically to the local mechanism implemented in the present computational framework. The mechanism allows the evolution of the parent fuel, oxygen, major combustion products, intermediate species, and radicals to be resolved throughout the engine cycle.

\paragraph{Reference operating conditions.}

The One-Zone thermo-chemical analysis is performed using the Mazda 13B
naturally aspirated Wankel engine as a fixed geometric reference. Keeping the geometry unchanged makes it possible to attribute subsequent changes in pressure, temperature, and chemical composition primarily to the
thermo-chemical formulation rather than to simultaneous changes in chamber geometry. The main Wankel geometric relations follow the classical rotary-engine description reported by Yamamoto \cite{yamamoto1981}.

The principal operating and numerical parameters of the reference case are summarized in Table~\ref{tab:onezone_conditions}.

\begin{table}[htbp]
\centering
\caption{Reference conditions of the One-Zone thermo-chemical model.}
\label{tab:onezone_conditions}
\begin{tabular}{lcc}
\hline
Parameter & Value & Unit \\
\hline
Eccentricity, $e$ & 15.0 & mm \\
Generating radius, $R$ & 105.0 & mm \\
Rotor width, $b$ & 80.0 & mm \\
Compression ratio & 9.4:1 & -- \\
Eccentric-shaft speed & 2750 & rpm \\
Equivalence ratio, $\phi$ & 1.0 & -- \\
Fresh-mixture AFR & 15.0276 & kg/kg \\
Intake temperature & 330 & K \\
Intake pressure & $0.95P_{\mathrm{atm}}$ & -- \\
Residual temperature & 650 & K \\
Residual pressure & $0.20P_{\mathrm{atm}}$ & -- \\
Exhaust temperature & 700 & K \\
Exhaust pressure & $P_{\mathrm{atm}}$ & -- \\
Wall temperature & 450 & K \\
IVC & 270 & deg \\
Ignition start & 525 & deg \\
Minimum-volume position & 540 & deg \\
Exhaust opening & 810 & deg \\
Chemical species & 56 & -- \\
Chemical reactions & 332 & -- \\
\hline
\end{tabular}
\end{table}

The stoichiometric air--fuel ratio is calculated from the species and elemental composition defined by the kinetic mechanism rather than being imposed as a fixed generic gasoline value.

\paragraph{Mass exchange.}
\paragraph{Mass exchange.}

Intake and exhaust processes are represented through upstream and downstream reservoirs connected to the variable-volume reactor during their corresponding angular windows. Mass and enthalpy fluxes are therefore included directly in the reactor evolution, allowing the One-Zone formulation to represent an open engine cycle rather than an isolated constant-mass reactor \cite{Cantera2025}.

During the closed interval from intake-valve closing (IVC) to exhaust opening (EO), the modeled chamber is treated as a closed system, providing an independent interval for first-law verification.

\paragraph{Effective ignition treatment.}

The current spark-ignition treatment uses a finite-duration thermal activation pulse beginning at

\begin{equation}
\theta_{\mathrm{spark}}=525^\circ
\end{equation}

and extending over $6^\circ$ of eccentric-shaft rotation. The pulse raises the reacting mixture toward a target ignition temperature of 1150 K, after which the evolution of the chemical species is governed by the finite-rate kinetic mechanism.

The associated thermal input is explicitly recorded in the first-law balance. This treatment constitutes a numerical ignition strategy for the present zero-dimensional reactor and should not be interpreted as the electrical energy delivered by a physical spark-discharge system.

\paragraph{Wall heat transfer.}

Heat transfer between the reacting gas and the engine housing is included during thermo-chemical integration according to

\begin{equation}
\dot{Q}_{\mathrm{wall}}
=
hA_w\left(T-T_{\mathrm{wall}}\right),
\end{equation}

where $A_w$ is the equivalent instantaneous wetted area and
$T_{\mathrm{wall}}=450$ K.

The convective heat-transfer coefficient is estimated using a reduced-order Woschni-type correlation. Woschni originally proposed an empirical formulation for estimating instantaneous convective heat transfer in internal combustion engines \cite{Woschni1967}. The adapted expression implemented in the present model is

\begin{equation}
h =
0.685
\left(\frac{N}{1000}\right)^{0.786}
\left(\frac{P}{10^5}\right)^{0.786}
T^{-0.525}
\times10^3,
\end{equation}

where $N$ is the eccentric-shaft speed in rpm, $P$ is the instantaneous chamber pressure in Pa, and $T$ is the gas temperature in K.

Because both the correlation and the equivalent wetted-area calculation are reduced-order approximations, the present treatment should be interpreted as a Woschni-type heat-transfer model rather than as an exact Wankel-specific wall-heat-transfer formulation.

\paragraph{Numerical integration and multi-cycle convergence.}

The reacting system is integrated using the adaptive reactor-network solver provided by Cantera \cite{Cantera2025}. A general output resolution of $0.5^\circ$ of eccentric-shaft rotation is employed throughout the cycle, while the output grid is refined to $0.05^\circ$ within the combustion region. The maximum permitted internal integration step corresponds to $0.25^\circ$, with relative and absolute solver tolerances of

\begin{equation}
RTOL=10^{-9},
\qquad
ATOL=10^{-15}.
\end{equation}

Successive engine cycles are calculated until periodic operation is achieved. Pressure, temperature, total chamber mass, and chemical composition are compared between consecutive cycles. The periodicity criterion requires the corresponding closure errors to remain below 1\%, with a minimum of three simulated cycles before convergence is accepted.

This procedure reduces the dependence of the reported thermo-chemical state on the arbitrary initial reactor conditions and establishes a periodic reference state for the subsequent thermodynamic analysis.

\paragraph{Fuel conversion and conversion-based combustion phasing.}

The disappearance of the parent fuel species is used as a chemical progress variable. The absolute parent-fuel conversion is defined relative to the fuel mass present at intake-valve closing,

\begin{equation}
X_f(\theta_e)
=
\frac{m_{f,\mathrm{IVC}}-m_f(\theta_e)}
{m_{f,\mathrm{IVC}}}.
\end{equation}

Conversion-based phasing markers CA10, CA50, and CA90 are defined as the eccentric-shaft angles at which $X_f$ reaches 10\%, 50\%, and 90\%, respectively.

These quantities are obtained directly from the finite-rate disappearance of the parent fuel species and are not imposed through a Wiebe function. Consequently, the CA10, CA50, and CA90 values reported in the present One-Zone model should be interpreted specifically as fuel-conversion-based markers rather than experimental mass-fraction-burned measurements.

\paragraph{Thermodynamic reconstruction and energy conservation.}

The thermodynamic post-processing stage uses the complete multi-species state obtained from the reacting system to evaluate the instantaneous properties of the working mixture. Consequently, properties such as internal energy, enthalpy, and specific heats evolve with temperature and chemical composition.

The indicated work is evaluated directly from the pressure--volume trajectory,

\begin{equation}
W=\oint P\,dV.
\end{equation}

For the complete open cycle, the first-law balance is written as

\begin{equation}
\Delta U
=
H_{\mathrm{in}}
-
H_{\mathrm{out}}
+
Q_{\mathrm{ign}}
-
Q_{\mathrm{wall}}
-
W,
\end{equation}

consistent with the general control-volume energy formulation
\cite{cengel2015thermodynamics}.

A second and more restrictive verification is performed over the closed interval between intake-valve closing and exhaust opening,

\begin{equation}
\Delta U_{\mathrm{IVC}\rightarrow\mathrm{EO}}
=
Q_{\mathrm{ign}}
-
Q_{\mathrm{wall}}
-
W_{\mathrm{IVC}\rightarrow\mathrm{EO}}.
\end{equation}

Chemical reaction energy is already contained in the internal energy and enthalpy of the reacting multi-species mixture. Therefore, the diagnostic chemical heat-release rate obtained from the reaction rates is not introduced again as an independent heat source in the first-law balance, avoiding double counting of chemical energy.
 
\subsubsection{2.3.2 two-Zone Model}
\subsubsection{2.3.3 Two-Zone Gasoline-Oil Model}

\subsection{2.4 Torque and Power Analysis via Vectorial Mechanics}

\subsubsection{2.4.1 Vectorial Position Analysis of the Rotor}

To complement the thermodynamic and geometric modeling, a vector-based mechanical analysis approach was implemented to analyze the position of key points on the rotor throughout the engine cycle. This method provides a dynamic view of the forces and torques generated by combustion and how they act on the rotor.



The position of the rotor's centroid \( \vec{C}_r \) at each angular step \( \theta \) is given by:
\begin{equation}
\vec{C}_r(\theta) = [e \cdot \cos(\theta),\ e \cdot \sin(\theta)]
\end{equation}
where \( e \) is the eccentricity of the rotor's path.


The angular position of the rotor itself is governed by the relation \(
\alpha(\theta) = \frac{1}{3}\theta
\), which ensures proper phasing relative to the housing epitrochoid. At each step, the rotor triangle is rotated and translated, and the center of one of its sides (e.g., side AB) is tracked.

Two vectors are calculated for each frame:
\begin{itemize}
    \item \( \vec{r}_1 \): from the rotor centroid to the center of side AB.
    \item \( \vec{r}_2 \): from the engine axis (origin) to the rotor centroid.
\end{itemize}

The total position vector of the point of application of force is:
\begin{equation}
\vec{r}_{\text{total}} = \vec{r}_1 + \vec{r}_2,
\end{equation}
providing the arm for torque calculation. As shown in the figure


\begin{figure}[H] % or [htbp]
  \centering
  \includegraphics[width=0.5\textwidth]{Images/Torqueexplicado.jpg}
  \caption{Representation of the previously mentioned vectors in the wankel engine.}
  \label{fig:torque_expl}
\end{figure}


The direction of the applied force is considered opposite to \( \vec{r}_1 \), representing the pressure exerted on the rotor face during expansion. This direction is normalized to obtain a unit vector \( \hat{f} \), and the instantaneous torque is then calculated as a cross product:
\begin{equation}
\tau(\theta) = \vec{r}_{\text{total}}(\theta) \times \left(F(\theta) \cdot \hat{f}(\theta)\right),
\end{equation}
where \( F(\theta) \) is the force derived from pressure multiplied by the chamber area at that moment.

This approach enables a dynamic, time-resolved calculation of torque, revealing asymmetries and time dependencies within the engine cycle that are not visible from thermodynamic equations alone. The methodology bridges geometric modeling and dynamic vector analysis, improving the accuracy of performance simulations and enabling optimization based on vectorial torque profiles.


To analyze the mechanical response, the torque \( \tau \) is computed using the vector cross product between the arm of force \( \vec{r} \) (from the origin to the point of force application) and the force vector \( \vec{F} \):
\begin{equation}
\tau = \vec{r} \times \vec{F} = r_x F_y - r_y F_x
\end{equation}
The direction of the force is assumed opposite to the vector connecting the rotor center to the point of application. The torque is computed at each time step using the pressure-based force and the position vector of the rotor.

Power output is computed as:
\begin{equation}
P = \tau \cdot \omega
\end{equation}

and the net mechanical work is estimated from:
\begin{equation}
W_{\text{net}} = \oint P \, dV
\end{equation}

This analysis helps characterize both instantaneous and average performance of the engine.


\section{3. Results}


\subsection{3.1 Geometrical Validation of the Rotor and Chamber}
 The geometry of the rotor and housing was first implemented using parametric equations derived from the epitrochoidal path and the kinematic behavior of the Wankel mechanism. The resulting rotor shape and apex trajectory were validated graphically, confirming their periodic and symmetric nature. This validation was further supported by simulating the rotor’s motion and comparing the results with the analytical representation obtained from Equations (1) and (2). This step provided a reliable geometric foundation for the subsequent volume and thermodynamic analyses.

\begin{figure}[H] % or [htbp]
  \centering
  \includegraphics[width=1\textwidth]{Images/comparacion_trayectorias.eps}
  \caption{Comparison between the theoretical epitrochoid (left), obtained from the analytical formulation given by Equations (1) and (2), using parámeters R = 70.71 and e = 10.06 with the side of the triangle for the simulation L = R$\sqrt{3}$ and the simulated vertex trajectory (right). The close match between both curves validates the geometric implementation of the rotor’s motion and confirms the accuracy of the kinematic model.}
  \label{fig:torque_expl}
\end{figure}
\subsection{3.2 Geometrical Simulation}

By taking Equations (5), (6), and (7), it is possible to construct the graphical representation of the three combustion chambers as a function of the angular position of the eccentric shaft. Each equation defines the instantaneous volume of one chamber as a function of the rotation angle. The chambers are geometrically identical but angularly phased by 2$\pi$/3 radians (120°). 

\begin{figure}[H] % or [htbp]
  \centering
  \includegraphics[width=1\textwidth]{Images/volumes_vs_theta.eps}
  \caption{Volume of the engine chambers with respect to the rotor angle.}
  \label{fig:torque_expl}
\end{figure}

By varying key geometric parameters of the rotary Wankel engine, such as the rotor radius R and the eccentricity e, it is possible to obtain different rotor–housing configurations that directly influence the shape of the combustion chambers. These geometric changes modify the maximum and minimum volumes reached during the engine cycle, thereby altering the compression ratio. Increasing e, for example, typically reduces the minimum chamber volume, leading to a higher compression ratio, while changing R can affect both the displacement and the overall chamber geometry. As a result, adjusting these parameters allows for the design of engines with tailored performance characteristics, enabling optimization for power output, efficiency, or specific operating conditions.

In this study, three different geometries will be compared, starting from Geometry 1 as the baseline. The variations will be introduced by modifying the parameters R and e, allowing the influence of these geometric changes on the engine’s performance to be analyzed

\begin{figure}[H] % or [htbp]
  \centering
  \includegraphics[width=0.6\textwidth]{Images/164.eps}
  \caption{Geometry 1 with a compression ratio of 1:6.4, with parameters R = 94.04 mm, e = 20 mm, and W = 30 mm.}
  \label{fig:torque_expl}
\end{figure}

\begin{figure}[H] % or [htbp]
  \centering
  \includegraphics[width=0.6\textwidth]{Images/183.eps}
  \caption{Geometry 2 with a compression ratio of 1:8.3, with parameters R = 100 mm, e = 24.1 mm, and W = 30 mm.}
  \label{fig:torque_expl}
\end{figure}

\begin{figure}[H] % or [htbp]
  \centering
  \includegraphics[width=0.6\textwidth]{Images/196.eps}
  \caption{Geometry 3 with a compression ratio of 1:9.6, with parameters R = 94.04 mm, e = 24.1 mm, and W = 30 mm.}
  \label{fig:torque_expl}
\end{figure}

The following results present the comparison of three different geometries of the rotary chamber. Pressure and temperature profiles are shown as functions of the crank angle for each configuration, allowing us to observe how the geometric parameters influence the thermodynamic behavior during compression and expansion. In addition, a P–V diagram is provided for the three geometries using gasoline as the fuel, which highlights the differences in the enclosed cycle area and therefore the work output of each case.

\begin{figure}[H] % or [htbp]
  \centering
  \includegraphics[width=1\textwidth]{Images/pressure.eps}
  \caption{Pressure plot for the three geometries.}
  \label{fig:torque_expl}
\end{figure}

The pressure curves for the three geometries highlight the strong influence of the eccentricity e and the radius R on the peak values reached during the cycle. Geometry G1 (R = 94.04 mm, e = 20 mm, W = 30 mm) shows the lowest maximum pressure, indicating a lower compression effect compared to the other two cases. Geometry G2 (R = 100 mm, e = 24.1 mm, W = 30 mm) achieves a higher peak pressure due to the larger chamber radius combined with a greater eccentricity, which increases the compression ratio. Geometry G3 (R = 94.04 mm, e = 24.1 mm, W = 30 mm) reaches the highest peak pressure of all, showing that eccentricity is a critical parameter in enhancing the compression and the resulting pressure rise after combustion.

Despite these differences in maximum pressure, all geometries show a similar qualitative trend: a gradual pressure increase during compression, a sudden rise due to isochoric heat addition, and a subsequent decrease during expansion. The magnitude of the peak pressure directly affects the indicated work per cycle, suggesting that G3 would deliver the highest power output, followed by G2, while G1 would provide the lowest.


\begin{figure}[H] % or [htbp]
  \centering
  \includegraphics[width=1\textwidth]{Images/temperature.eps}
  \caption{Temperature plot for the three geometries.}
  \label{fig:torque_expl}
\end{figure}

The temperature profiles for the three geometries show the typical evolution of the cycle: a moderate rise during compression, a sharp jump after heat release, and a smooth decrease during expansion. The differences among the geometries are mainly observed in the magnitude of the peak temperature.

Geometry G3 (R = 94.04 mm, e = 24.1 mm, W = 30 mm) produces the highest maximum temperature. This is because the lower eccentricity reduces the effective chamber volume at compression, increasing the temperature rise when heat is released.

Geometry G2 (R = 100 mm, e = 24.1 mm, W = 30 mm) shows slightly lower peak temperatures than G1. The larger rotor radius combined with higher eccentricity increases the chamber volume, which spreads the added heat over more mass, lowering the temperature peak.

Geometry G1 (R = 94.04 mm, e = 20 mm, W = 30 mm) reaches the lowest maximum temperature of the three. The combination of smaller radius with greater eccentricity further enlarges the compressed volume, reducing the intensity of the temperature rise.

Overall, while the three geometries exhibit the same qualitative behavior, peak temperature decreases as eccentricity and chamber volume increase, highlighting how geometric design strongly influences the thermal conditions of the working fluid.

\begin{figure}[H] % or [htbp]
  \centering
  \includegraphics[width=1\textwidth]{Images/pv_diagram.eps}
  \caption{P-V plot for the three geometries.}
  \label{fig:torque_expl}
\end{figure}

The P–V diagram clearly shows how geometry affects the thermodynamic cycle of the rotary chamber. While all three configurations follow the same qualitative behavior—compression, isochoric heat addition, and expansion—the enclosed cycle areas and peak pressures differ significantly.

For \textbf{G1} (\(R=94.04\ \text{mm}\), \(e=20\ \text{mm}\), \(W=30\ \text{mm}\)), the cycle exhibits the lowest work, \(W_{\mathrm{net}}=0.737\ \text{kJ}\), together with the highest specific fuel consumption, \(\mathrm{SFC}=0.038\ \text{g/kJ}\). The corresponding peak values are \(P_{\max}=8158.78\ \text{kPa}\) and \(T_{\max}=3521.67\ \text{K}\). In \textbf{G2} (\(R=100\ \text{mm}\), \(e=24.1\ \text{mm}\), \(W=30\ \text{mm}\)), the net work rises to \(1.024\ \text{kJ}\) (\(+38.9\%\) relative to G1), while SFC decreases to \(0.033\ \text{g/kJ}\) (\(-13.2\%\)). Peak pressure reaches \(10742.99\ \text{kPa}\) (\(+31.7\%\)) with \(T_{\max}=3570.52\ \text{K}\) (\(+1.39\%\)). Finally, \textbf{G3} (\(R=94.04\ \text{mm}\), \(e=24.1\ \text{mm}\), \(W=30\ \text{mm}\)) sustains high work, \(1.005\ \text{kJ}\) (\(+36.4\%\) vs G1), and achieves the lowest SFC, \(0.031\ \text{g/kJ}\) (\(-18.4\%\)), while exhibiting the highest peak pressure, \(P_{\max}=12467.94\ \text{kPa}\) (\(+52.8\%\) vs G1), and \(T_{\max}=3598.84\ \text{K}\) (\(+2.19\%\)).

Comparatively, increasing eccentricity from \(e=20\) to \(e=24.1\) at constant \(R\) (G1 \(\rightarrow\) G3) is associated with simultaneous gains in \(W_{\mathrm{net}}\) and reductions in SFC, at the expense of a pronounced rise in \(P_{\max}\), while \(T_{\max}\) changes little (within \(\leq 2.2\%\)). Increasing radius from \(R=94.04\ \text{mm}\) to \(R=100\ \text{mm}\) at constant \(e\) (G3 \(\rightarrow\) G2) moderates \(P_{\max}\) by about \(14\%\) relative to G3, maintains a similarly high \(W_{\mathrm{net}}\), and incurs only a slight SFC penalty.



Overall, the comparison demonstrates that increasing eccentricity (e) raises the compression ratio and cycle efficiency, while increasing radius (R) enlarges volume but slightly reduces the maximum pressure. Thus, G3 offers the most favorable balance for maximizing indicated work output, followed by G2, with G1 being the least effective.

\subsection{3.3 Baseline Ideal-Cycle Comparison of Different Fuels} 
The results presented in this subsection correspond to the simplified
ideal-cycle formulation introduced in Section~2.2. They are retained as a baseline fuel-sensitivity study and should not be interpreted as a detailed finite-rate chemical comparison. The thermo-chemical results obtained from the kinetic model are presented separately in the following subsection.

The numerical model developed in this study provides a detailed representation of the thermodynamic behavior of the rotary Wankel engine. The results include the evolution of pressure during compression and expansion, as well as the corresponding temperature changes. These outputs were computed for a defined engine configuration, considering ambient conditions and the geometric parameters $R = 94.04 \ \text{mm}$, $e = 24.1 \ \text{mm}$, $W = 30 \ \text{mm}$, and a compression ratio of $1\!:\!9.6$. Furthermore, three different fuels were used to compare the thermodynamic cycle: gasoline, methane, and bioethanol.To develop this thermodynamic study, the lower heating values (LHV) of the fuels were considered: 
gasoline $41.1 \ \text{MJ}\,\text{kg}^{-1}$ \cite{ornl_tedb37}, methane $50.0 \ \text{MJ}\,\text{kg}^{-1}$ \cite{turns2013}, 
and bioethanol $26.8 \ \text{MJ}\,\text{kg}^{-1}$ \cite{demirbas2008}. 
These values were used to calculate the energy released per gram of fuel in the combustion chamber of the Wankel engine.


The following figures illustrate the behavior of pressure, temperature, and the P–V diagram for three different fuels – gasoline (iso-octane), methane ($CH_4$), and bioethanol ($C_2H_5OH$) – considering the previously defined rotor and housing geometry.

\begin{figure}[H] % or [htbp]
  \centering
  \includegraphics[width=1\textwidth]{Images/pressure_fuels.eps}
  \caption{Pressure during the compression and expansion phases of the engine for 3 different fuels.}
  \label{fig:torque_expl}
\end{figure}

The three pressure traces follow the same pattern: a gradual rise during compression, a sharp jump at the heat-release point, and a smooth decay through expansion. In this comparison, methane shows the highest peak pressure, gasoline is slightly lower, and bioethanol is the lowest. Methane’s peak is highest because, even though its stoichiometric mixture uses less fuel per unit air, the mixture’s effective heat capacity is lower and the heat release raises temperature—and thus pressure—more sharply. Gasoline lands close behind due to strong heating but a slightly “heavier” mixture that tempers the temperature rise. Bioethanol trails because its lower heating value and richer mixture (plus higher specific heat of the fuel) spread the added energy over a more heat-absorbent charge, yielding a smaller pressure spike. The overall curve shapes are similar across fuels; what changes most is the magnitude of the peak.

\begin{figure}[H] % or [htbp]
  \centering
  \includegraphics[width=1\textwidth]{Images/temperature_fuels.eps}
  \caption{Temperature during the compression and expansion phases of the engine for 3 different fuels.}
  \label{fig:torque_expl}
\end{figure}

The temperature traces for gasoline, methane, and bioethanol share the same shape: a steady rise during compression, a sharp step at the isochoric heat-release point (near $\theta$ $\approx$ $\pi$/2), and a monotonic drop through expansion. The peak ordering is Gasoline > Methane > Bioethanol. Gasoline attains the highest maximum because, for the same trapped air mass, its combination of heating value and mixture properties yields the largest effective sensible heating at constant volume. Methane peaks slightly lower—its leaner stoichiometric mixture means less fuel mass per unit air, moderating the temperature jump. Bioethanol shows the lowest peak, consistent with its lower LHV and higher specific heat, which increase the mixture’s heat-absorbing capacity and limit the rise.

After the peak, all three curves cool at comparable rates as the volume increases and the gases expand adiabatically. The differences among fuels are therefore most visible at and immediately after the heat-addition event; away from that region, the curves converge, underscoring that fuel choice primarily shifts the magnitude of peak temperature rather than the overall trend.

\begin{figure}[H] 
  \centering
  \includegraphics[width=1\textwidth]{Images/pv3c.eps}
  \caption{P-V diagram of the cycle with 3 different fuels.}
  \label{fig:torque_expl}
\end{figure}

The P--V diagram compares the thermodynamic performance of gasoline (iso-octane), methane (CH$_4$), and bioethanol (C$_2$H$_5$OH) under identical geometric conditions. While the overall cycle shapes remain similar, the peak pressures, maximum temperatures, and specific fuel consumptions highlight the distinct characteristics of each fuel.  

Methane achieves the highest final pressure after the isochoric process (\textbf{11,699.19 kPa}) and a peak temperature of about \textbf{3637.88 K}. This leads to the largest work output per cycle (\textbf{0.865 kJ}) and the lowest specific fuel consumption (\textbf{0.031 g/kJ}). These results confirm methane’s superior thermodynamic efficiency in this configuration.  

Gasoline (iso-octane) exhibits a final pressure of \textbf{11,607.89 kPa} and a maximum temperature of \textbf{3609.49 K}. Its net work per cycle is slightly lower than methane at \textbf{0.859 kJ}, with a higher specific fuel consumption (\textbf{0.037 g/kJ}). This positions gasoline as the second-best performer, close to methane in terms of work but less favorable in fuel efficiency.  

Bioethanol displays the lowest final pressure (\textbf{10,623.15 kPa}) and temperature (\textbf{3383.28 K}) among the three fuels. Consequently, its net work per cycle (\textbf{0.799 kJ}) is the smallest, and its specific fuel consumption (\textbf{0.065 g/kJ}) is significantly higher. These results underscore the reduced energy density of bioethanol and its comparative disadvantage in efficiency under the same operating conditions.  
  

Because this study relies on a cycle-level thermodynamic simulation without a coupled load or dynamometer model, we cannot report SBHP (brake specific horsepower). SBHP is derived from brake power, which requires shaft torque under an externally applied load and must account for mechanical, friction, and pumping losses. Our model does not simulate the load path or include a validated loss map; therefore brake torque—and thus brake power and SBHP—cannot be resolved with sufficient fidelity. Any SBHP value would be speculative. Instead, we report indicated quantities (e.g., IMEP, indicated work/power) and theoretical torque; obtaining SBHP would require experimental dynamometer data or an extended model that explicitly includes a load and calibrated mechanical losses.


Overall, the diagram and data confirm that methane provides the best performance in terms of pressure, work output, and fuel efficiency, followed closely by gasoline, while bioethanol, although cleaner and renewable, delivers less energy per cycle and requires more fuel mass for the same output.

\subsection{3.4 Vectorial Torque Analysis}

Torque was computed from the normal force on the active rotor side and the lever arm with respect to the origin,
\begin{equation}
  \boldsymbol{\tau} = \mathbf{r} \times \mathbf{F}, \qquad
  \mathbf{F} = \bigl(P(\theta)-P_{\mathrm{atm}}\bigr)\,A\, \qquad
  A = L\,W,
\end{equation}
where $A$ is the effective thrust area in m$^2$, $L=\sqrt{3}\,R$ is the rotor side length, $W$ is the axial width. The pressure trace $P(\theta)$ comes from the thermodynamic cycle model (adiabatic compression, isochoric heat addition, adiabatic expansion). 

\subsection{Torque vs. angle for three geometries}

 
Figure~\ref{fig:tau_geoms} compares three geometries under the same fuel:

\begin{figure}[H]
   \centering
   \includegraphics[width=1\linewidth]{Images/torque.eps}
   \caption{Torque vs.\ angle for the three geometries.}
   \label{fig:tau_geoms}
\end{figure} 

We compare three geometries under identical operating and fuel assumptions:
G1 $(R{=}94.04,\ e{=}20,\ W{=}30)$,
G2 $(R{=}100,\ e{=}24.1,\ W{=}30)$, and
G3 $(R{=}94.04,\ e{=}24.1,\ W{=}30)$.

Peaks occur just after the isochoric step, near $\theta \approx \pi/2$:
\begin{table}[h]
  \centering
  \begin{tabular}{lcccc}
    \toprule
    Geometry & Peak torque $\tau_{\max}$ [N$\cdot$m] & Angle [rad] & vs.\ G1 \\
    \midrule
    G1 & 127.466 & 1.59 & --- \\
    G2 & 215.482 & 1.59 & +69.1\% \\
    G3 & 235.381 & 1.59 & +84.7\% \\
    \bottomrule
  \end{tabular}
\end{table}

Relative peak comparisons: G1$\to$G2 $+69.1\%$, G1$\to$G3 $+84.7\%$, G2$\to$G3 $+9.2\%$.


\subsection{Torque vs.\ angle for three fuels (fixed geometry)}

\begin{figure}[H]
   \centering
   \includegraphics[width=1\linewidth]{Images/torque_fuels.eps}
   \caption{Torque vs.\ angle for the three different fuels with the same geometry.}
   \label{fig:tau_geoms}
\end{figure} 


\label{sec:tau_fuels}
Geometry held constant at  $R = 94.04 \ \text{mm}$, $e = 24.1 \ \text{mm}$, $W = 30 \ \text{mm}$, and a compression ratio of $1\!:\!9.6$. 
We compare \emph{Gasoline (iso-octane)}, \emph{Methane (CH$_4$)}, and \emph{Bioethanol (C$_2$H$_5$OH)} under identical cycle assumptions (adiabatic compression, isochoric heat addition, adiabatic expansion).



% Requires \usepackage{booktabs} in the preamble

\subsection*{Peak torque}
Maximum torque from the torque--angle curves:

\begin{table}[h]
  \centering
  \begin{tabular}{lcccc}
    \toprule
    Fuel & $\tau_{\max}$ [N$\cdot$m] & Angle [rad] & vs.\ Gasoline & vs.\ Methane \\
    \midrule
    Gasoline (iso-octane)      & 228.032 & 1.59 & ---     & $-3.8\%$ \\
    Methane (CH$_4$)           & 237.056 & 1.59 & $+4.0\%$ & --- \\
    Bioethanol (C$_2$H$_5$OH)  & 227.136 & 1.59 & $-0.4\%$ & $-4.2\%$ \\
    \bottomrule
  \end{tabular}
  
  \caption{Comparison of peak torque ($\tau_{\max}$) and peak angle for each fuel. Relative columns report percentage difference in $\tau_{\max}$.}
  \label{tab:max_tau_fuels}
\end{table}



All fuels reach their maximum torque at essentially the same crank angle ($\theta \approx 90.72^\circ$), just after the isochoric heat addition. The minimum torque occurs slightly before that ($\theta \approx 89.62^\circ$) and is similar for all fuels ($\tau_{\min}\approx -46.27$ N$\cdot$m). In terms of amplitude, methane shows the highest peak (237.06 N$\cdot$m), about $+4.0\%$ above gasoline (228.03 N$\cdot$m) and $+4.2\%$ above bioethanol (227.14 N$\cdot$m). The negative torque magnitude during compression is around $20\%$ of each fuel’s peak, indicating the same basic pulse shape across fuels; the fuel choice mainly scales the torque level rather than shifting its phase.

\subsection{Finite-Rate Thermo-Chemical Results}

The following results correspond to the current One-Zone thermo-chemical baseline using iso-octane as the reference hydrocarbon fuel. This model constitutes the first numerically verified stage of the thermo-chemical development and will subsequently be used as the reference case for comparison with the Two-Zone formulation and, ultimately, with the Two-Zone gasoline-lubricating-oil model.

\subsubsection{Periodic Multi-Cycle Convergence}

The One-Zone iso-octane case satisfies the prescribed periodicity criterion after three complete engine cycles. At the end of the third cycle, the closure errors are approximately $2.06\times10^{-5}\%$ for pressure, $0.0117\%$ for temperature, $0.0101\%$ for chamber mass, and
$7.12\times10^{-5}\%$ for chemical composition.

All four closure quantities remain well below the adopted 1\% periodicity criterion. The third cycle is therefore adopted as the reference periodic solution for the subsequent chemical, thermodynamic, and mechanical analyses.

\begin{figure}[H]
  \centering
  \includegraphics[width=1\textwidth]{Images/onezone_multicycle_convergence.png}
  \caption{Multi-cycle convergence of the One-Zone iso-octane model.}
  \label{fig:onezone_multicycle}
\end{figure}

\subsubsection{Pressure and Temperature Evolution}

The converged One-Zone cycle reaches a maximum chamber pressure of
approximately $5.772$ MPa and a maximum gas temperature of $3058.8$ K.
The rapid increase in both variables is associated with the main
fuel-conversion interval near the minimum-volume region, after which pressure and temperature decrease progressively during expansion.

Unlike the baseline ideal-cycle formulation, these profiles result from the coupled evolution of variable chamber volume, finite-rate chemical kinetics, mass exchange, effective ignition, and wall heat transfer.

\begin{figure}[H]
  \centering

  \begin{minipage}[b]{0.49\textwidth}
    \centering
    \includegraphics[width=\textwidth]{Images/onezone_pressure_iso_octane.png}
    \small (a) Chamber pressure.
  \end{minipage}
  \hfill
  \begin{minipage}[b]{0.49\textwidth}
    \centering
    \includegraphics[width=\textwidth]{Images/onezone_temperature_iso_octane.png}
    \small (b) Chamber temperature.
  \end{minipage}

  \caption{Pressure and temperature evolution as a function of eccentric-shaft angle for the converged One-Zone iso-octane thermo-chemical simulation.}
  \label{fig:onezone_pressure_temperature}
\end{figure}
Unlike the ideal-cycle results presented previously, these profiles emerge from the coupled solution of variable chamber volume, chemical kinetics, mass exchange, effective ignition, and wall heat transfer.

\subsubsection{Fuel Conversion and Combustion Phasing}

The parent iso-octane species is fully consumed before exhaust opening in the converged One-Zone solution. This result corresponds to 100\% conversion of the tracked parent fuel species and should not, by itself, be interpreted as complete oxidation of all intermediate hydrocarbon species or complete
combustion of a real gasoline fuel.

The conversion-based phasing markers obtained from the finite-rate solution are

\begin{equation}
CA10=530.23^\circ,
\qquad
CA50=531.16^\circ,
\qquad
CA90=531.85^\circ.
\end{equation}

Since the effective ignition event begins at $525^\circ$ and the
minimum-volume position occurs at $540^\circ$, the main parent-fuel conversion interval is concentrated between ignition and the minimum-volume region. The interval between CA10 and CA90 is

\begin{equation}
\Delta CA_{10-90}=1.62^\circ.
\end{equation}

\begin{figure}[H]
  \centering
  \includegraphics[width=1\textwidth]{Images/onezone_iso_octane_conversion.png}
  \caption{Absolute iso-octane conversion as a function of eccentric-shaft angle. The locations of the spark event, CA10, CA50, and CA90 are obtained directly from the finite-rate chemical solution.}
  \label{fig:onezone_conversion}
\end{figure}

\subsubsection{Evolution of the Main Chemical Species}

The finite-rate formulation provides the evolution of the chemical composition throughout the complete engine cycle. For the present iso-octane mechanism, the principal species selected for visualization are
$\mathrm{IC_8H_{18}}$, $\mathrm{O_2}$, $\mathrm{CO_2}$,
$\mathrm{H_2O}$, CO, $\mathrm{H_2}$, and $\mathrm{N_2}$.

The disappearance of the parent fuel and oxygen during the main reaction interval is accompanied by the formation of the principal combustion products. At exhaust opening, the calculated parent iso-octane conversion is 100\% and the oxygen conversion is approximately 99.88\%. The calculated $\mathrm{CO_2}$ and $\mathrm{H_2O}$ masses correspond to approximately 97.96\% and 99.96\%, respectively, of the theoretical values obtained from elemental balance.

These species histories provide a chemical description that is not available from the baseline ideal heat-addition model. However, they remain predictions of the current reduced One-Zone kinetic formulation and should not yet be interpreted as experimentally validated engine-out emissions.

\begin{figure}[H]
  \centering
  \includegraphics[width=1\textwidth]{Images/onezone_main_species_iso_octane.png}
  \caption{Evolution of the main chemical species during the complete One-Zone iso-octane/air engine cycle.}
  \label{fig:onezone_species}
\end{figure}

\subsubsection{Indicated Work and Pressure-Volume Cycle}

The pressure-volume trajectory calculated from the converged thermo-chemical state produces a net indicated work of

\begin{equation}
W_{\mathrm{net}}=595.70~\mathrm{J}
\end{equation}

per modeled chamber cycle. The corresponding net indicated mean effective pressure is

\begin{equation}
IMEP_{\mathrm{net}}=909.85~\mathrm{kPa}.
\end{equation}

The closed-cycle gross indicated work is approximately $625.73$ J, while the calculated pumping contribution is approximately $-30.03$ J. The resulting pressure-volume trajectory provides the direct link between the thermo-chemical solution and the subsequent vectorial mechanical analysis.

\begin{figure}[H]
  \centering
  \includegraphics[width=1\textwidth]{Images/onezone_pv_iso_octane.png}
  \caption{Pressure--volume diagram obtained from the converged One-Zone iso-octane thermo-chemical simulation.}
  \label{fig:onezone_pv}
\end{figure}

\subsubsection{Energy-Balance Verification}

The integrated wall heat transfer over the complete cycle corresponds to an  energy loss of approximately

\begin{equation}
Q_{\mathrm{wall}}=166.22~\mathrm{J}.
\end{equation}

The principal first-law verification is performed over the closed interval between intake-valve closing and exhaust opening. The final energy residual is

\begin{equation}
R_E=0.0199~\mathrm{J}.
\end{equation}

Relative to the chemical energy associated with the trapped parent fuel, this corresponds to

\begin{equation}
\frac{R_E}{m_f LHV}\times100
=
1.52\times10^{-3}\%.
\end{equation}

The small residual indicates a consistent numerical closure between the change in internal energy, effective ignition input, wall heat transfer, and pressure-volume work over the closed interval.

\begin{figure}[H]
  \centering

  \begin{minipage}[b]{0.49\textwidth}
    \centering
    \includegraphics[width=\textwidth]{Images/onezone_first_law_balance.png}
    \small (a) First-law energy balance.
  \end{minipage}
  \hfill
  \begin{minipage}[b]{0.49\textwidth}
    \centering
    \includegraphics[width=\textwidth]{Images/onezone_energy_residual.png}
    \small (b) Energy-balance residual.
  \end{minipage}

  \caption{First-law validation over the closed interval from intake-valve closing to exhaust opening in the One-Zone iso-octane model.}
  \label{fig:onezone_energy_validation}
\end{figure}

The results obtained with the One-Zone formulation establish a numerically consistent thermo-chemical baseline for pressure, temperature, parent-fuel conversion, species evolution, indicated work, and energy conservation. Nevertheless, the homogeneous-reactor assumption assigns a single temperature and chemical composition to the complete chamber and therefore cannot resolve the thermal and compositional differences between burned and unburned regions.

For this reason, the next stage of the model introduces a Two-Zone formulation while preserving the present geometric reference, operating conditions, kinetic framework, and numerical verification criteria. Once this formulation is established, the chemical description will be extended from the present iso-octane reference case toward a multicomponent gasoline surrogate and, ultimately, a gasoline--lubricating-oil mixture.

\subsection{3.5 Comparison with a real engine}

In order to validate the numerical model, a comparison is planned between the simulated engine and a commercially available Wankel rotary engine, the AIE 650S 120 BHP. According to the manufacturer’s datasheet, this single-rotor engine has a displacement of 650 cm$^3$, delivers up to 120 bhp (90 kW) of power at 8000 rpm, and develops a torque of approximately 80 lb$\cdot$ft. The engine operates with a compression ratio of 9.6:1, features a liquid-cooled system (SPARCS\textsuperscript{*}), and supports multiple fuel types including AVGAS, gasoline, and heavy fuels. Its specific fuel consumption ranges between 310 and 350 g/kWh, and the core engine weight is 28 kg \cite{AIE650S2016}.


For the purposes of the simulation, the compression ratio and displacement were matched to the values of the real engine to ensure consistency in the thermodynamic conditions and chamber geometry. By setting these parameters as boundary conditions, the simulated model replicates the fundamental operating characteristics of the 650S. This alignment allows the subsequent comparative analysis to focus on performance indicators such as pressure evolution, net work per cycle, and specific fuel consumption, while preserving the equivalence in compression and displacement between the virtual and real engines.  

To enable a like–for–like comparison with the reference single-rotor engine, the simulated
stator–rotor geometry was calibrated using
$R=94.04\,\mathrm{mm}$,
$e=24.1\,\mathrm{mm}$, and
$W=30.0\,\mathrm{mm}$.
With these parameters, the chamber-volume function $V(\theta;R,e,W)$ yields
the maximum and minimum volumes $V_{\max}$ and $V_{\min}$ such that
\[
\mathrm{CR}=\frac{V_{\max}}{V_{\min}}=9.6.
\]
Adopting the standard Wankel convention (single rotor, three working chambers),
the nominal engine displacement is
 $650~\mathrm{cc}$ class of the real engine. This geometric matching
ensures that subsequent differences in performance arise from thermodynamic
assumptions rather than gross geometric discrepancies.

\begin{figure}[H]
   \centering
   \includegraphics[width=0.8\linewidth]{Images/engine_power_vs_rpm_bhp.eps}
   \caption{Real Engine Power vs RPMS.}
   \label{fig:power_realvscomp}
\end{figure} 

The graph compares the computational model (red dashed line) with the real engine data of the AIE 650S 120 BHP (blue line) as a function of engine speed.

A key observation is that the slope of both curves is similar over most of the analyzed range. In both cases, power increases approximately proportionally with RPM, indicating that the computational model captures the general trend of how power scales with rotational speed. This suggests that the underlying thermodynamic formulation of the model is consistent with the fundamental relationship between torque and angular velocity.

However, the main difference lies in the cut-off point of each curve. The real engine curve begins to flatten as it approaches 8000 rpm, reaching a maximum power of about 120 bhp. This reflects the physical limitations of the engine, such as reduced volumetric efficiency at high speeds, increased mechanical losses, combustion time constraints, and thermal effects. In contrast, the computational model does not show a saturation or peak within the same range. Instead, it continues increasing linearly, reaching values significantly higher than the real engine at 8000 rpm.

This indicates that although the model reproduces the general slope of the power increase, it does not yet incorporate the limiting mechanisms that define the real engine’s maximum power output. Therefore, the difference between the two curves is not primarily in their rate of growth, but in the absence of a realistic performance ceiling in the computational model.

Overall, the comparison demonstrates that the model correctly captures the power-RPM proportionality but requires refinement to accurately predict the engine’s peak power and high-speed behavior.

\section{4. Conclusions}

This study presents an integrative multi-physics framework for the computational analysis of the Wankel rotary engine, combining geometric, kinematic, thermodynamic, and mechanical models into a unified methodology. The proposed approach enables a coherent transition from the spatial definition of the engine geometry to the prediction of dynamic performance parameters, providing a systematic tool for analyzing rotary engine operation and exploring design variations.

\subsection{4.1 Integrated Modeling Framework}

The robustness of the proposed methodology arises from the interaction of four complementary physical models that capture the main mechanisms governing the operation of the Wankel engine.

\subsubsection{4.1.1 Geometric Model}

The geometric model defines the physical boundaries of the combustion chambers through the epitrochoidal stator geometry. Using the parameters of eccentricity ($e$) and generating radius ($R$), the stator trajectory is described by Equations (1) and (2), while the rotor profile is established through Equations (3) and (4). This configuration guarantees continuous contact between the rotor apexes and the housing surface, which is essential for sealing and chamber formation. As a result, the model provides the fundamental basis for calculating instantaneous areas and volumes, enabling the determination of the engine’s geometric compression ratio.

\subsubsection{4.1.2 Kinematic Model}

The kinematic model describes the relative motion of the rotor and eccentric shaft. The angular relationship $\alpha(\theta)=\frac{1}{3}\theta$ ensures the appropriate phase relationship required for the rotary implementation of the four-stroke cycle. Based on this relationship, the instantaneous chamber volumes $V_1$, $V_2$, and $V_3$ (Equations 5–7) are derived, capturing the volumetric evolution of the three working chambers separated by a phase shift of $2\pi/3$ radians. This formulation provides the dynamic volumetric input required for the thermodynamic analysis.

\subsubsection{4.1.3 Thermodynamic Model}

The thermodynamic model represents the engine cycle using a modified Otto formulation adapted to rotary geometry. Compression and expansion processes are described using a polytropic relation (Equation 8), while combustion is modeled as an isochoric heat addition process. The resulting peak pressure ($P_3$) and temperature ($T_3$) are obtained through Equations (9) and (10), incorporating the thermophysical properties and heating value of the air–fuel mixture. This model enables the prediction of pressure evolution within the combustion chambers throughout the cycle.

\subsubsection{4.1.4 Mechanical Model}

The mechanical model translates thermodynamic pressure into mechanical output. Through a vector mechanics formulation, the rotor centroid position ($\vec{C}_r$) and the effective force arm ($\vec{r}_{\text{total}}$) are defined, allowing the torque generated by gas pressure to be computed via the cross-product formulation (Equation 13). From this relationship, instantaneous torque, work, and power output are obtained using Equations (14) and (15), providing a direct link between chamber pressure evolution and engine performance.

\subsection{4.2 Model Validation and Application Potential}

The predictive capability of the proposed framework was assessed by benchmarking the model against the commercial AIE 650S engine. By reproducing the engine displacement and compression ratio (9.6:1), the model successfully captured the proportional relationship between engine power and rotational speed observed in real operational data. This agreement supports the validity of the integrated modeling strategy for preliminary engine analysis.

Beyond validation, the framework demonstrated flexibility for evaluating alternative operating scenarios. Simulations with different fuels showed that methane produces a 52.8\% increase in peak pressure compared to the baseline case, while bioethanol leads to higher specific fuel consumption due to its lower energy density. These results highlight the capability of the model to support comparative analyses of fuel strategies and performance optimization.

\subsection{4.3 Future Research Directions}

Although the present framework provides a robust representation of the indicated thermodynamic cycle, several improvements can further enhance its predictive capability.

\begin{itemize}
    
    \item \textbf{Loss mechanisms:} Including frictional losses, pumping work, and heat transfer through the housing walls to enable the prediction of brake power and overall engine efficiency.
    
    \item \textbf{System-level integration:} Coupling the engine model with external load curves (e.g., propeller or dynamometer models) and exploring alternative rotor geometries beyond the classical epitrochoidal configuration to optimize power density for aeronautical or automotive applications.
\end{itemize}

Overall, the methodology presented in this work establishes a structured computational foundation for the analysis and conceptual design of Wankel rotary engines, facilitating future developments in performance optimization, alternative fuels, and advanced rotary engine architectures.



\bibliographystyle{apacite}
\bibliography{references} 
\end{document}